\chapter{Graduate Attributes, AI Use and Project Repository}

\section{Graduate attribute evidence}
\label{app:ga}

This section summarises how the work meets the graduate attributes assessed in this course (GA1, GA4, GA5, GA6, GA8 and GA9, and GA10 in part), with the sections that hold the evidence.

\subsection*{GA1: Problem solving}
The research questions were framed around measurable dynamic response rather than subjective comfort, and each was tied to a requirement and a verification method (Section~\ref{sec:requirements-mapping}). Design choices were reasoned from first principles: a 200\,Hz sampling rate from the ISO 2631-1 band and the sampling theorem, drift-free task timing and non-blocking bounded queues so that acquisition is never held up (Section~\ref{sec:dual-imu-architecture}), the acceleration magnitude as an orientation-tolerant signal (Section~\ref{sec:ride-characterisation}), speed normalisation because raw RMS mixes roughness with speed (Section~\ref{sec:speed-norm}), and speed from wheelbase over the front-to-rear delay (Section~\ref{sec:road-analysis}). Faults were root-caused from evidence and the causes changed the design (Section~\ref{sec:investigations}, Table~\ref{tab:investigations}): a flash-configuration boot loop isolated one setting at a time (Section~\ref{sec:inv-flash}); the Raspberry Pi~4 under-voltage traced to the modem's supply path, which gave the telemetry unit its own supply (Section~\ref{sec:inv-power}); a time-of-flight start-up in which one failed sensor stopped all three, redesigned to skip, time out and re-probe (Section~\ref{sec:inv-tof}); and the flash-write acquisition stall (Section~\ref{sec:local-logging}).

\subsection*{GA4: Investigations, experiments and data analysis}
Assumptions were treated as hypotheses and tested on the hardware with a test chosen to tell the candidate causes apart (Section~\ref{sec:investigations}): the GNSS reply format and the modem's start-up (Section~\ref{sec:inv-gnss}); the BMS service layout and checksum (Section~\ref{sec:inv-bms}); invalid records, separated from transfer damage by an indexed read-back (Section~\ref{sec:inv-json}); the camera's colour cast (Section~\ref{sec:inv-colour}); the Pi~4 under-voltage, matched to its cause by a 2\,s power-flag log (Section~\ref{sec:inv-power}); and the acquisition stall, confirmed by comparing logging on and off (Section~\ref{sec:local-logging}). The features were validated against a synthetic signal with known answers, the IMUs calibrated against gravity, and the camera-to-IMU offset bounded. The second road session was designed as an experiment (Section~\ref{sec:field-test-design}): counterbalanced speeds, repeated passes, ground truth independent of the system (tape, video, meters), criteria fixed before the data are seen, and an analysis with confidence intervals for repeatability (Section~\ref{sec:road-analysis}). The raw evidence is indexed in Appendix~\ref{app:evidence} and Section~\ref{sec:data-availability}.

\subsection*{GA5: Use of engineering tools}
PlatformIO, the Arduino framework and FreeRTOS were used for the firmware, with esptool and serial diagnostics for hardware verification, and host tests run before each flash (Appendix~\ref{app:devrecord}). Python tools were written for testing: scripted bench tests, a browser bench console, a road recorder, and an analysis package with its own tests (Sections~\ref{sec:bench-tools} and~\ref{sec:road-analysis}). LiteRT and Picamera2 run the detector, with per-stage latency measured from the sensor's frame timestamps (Section~\ref{sec:hazard-detection}). Enclosures were modelled in SOLIDWORKS through scripts and in OpenSCAD and printed (Appendix~\ref{app:mechanical}). The limits of each tool were checked rather than assumed, for example a BMS checksum formula that no vendor document confirmed and that proved wrong on the vehicle (Section~\ref{sec:bms-interface}), and the ISO~2631-1 weighting, whose check is stated as internal because no copy of the standard was at hand (Section~\ref{sec:road-analysis}).

\subsection*{GA6: Professional and technical communication}
The work is communicated to several audiences, each in its own form. This report is written for examiners, in plain prose with a stated word count, and was revised with the supervisor through regular meetings, a draft Introduction and Methodology, and a check against his guidance lectures and example theses. A 10-minute presentation with a slide deck, timed talking points and prepared answers is for the panel on 16 and 17 November, and an A1 poster and a live demonstration plan are for the Open Day on 18 November (\texttt{Presentation/}). Video evidence of the road sessions, which the supervisor agreed should be included, is prepared with faces and the number plate blurred (Appendix~\ref{app:road1}). For anyone who wants to check a number, the data index lists every file behind the report with its checksum (Section~\ref{sec:data-availability}), and the lab notebook records each day's actions as they happened.

\subsection*{GA8: Individual work}
The hardware bring-up, firmware, bench tools, analysis and report were done individually, with the supervisor consulted on specific questions. Components were selected, compared against the requirements and ordered by the author, for example the 8$\times$8-zone time-of-flight sensor chosen over a single-point one to separate left, centre and right (Section~\ref{sec:tof-sensor}); each was tested on the bench before being relied on; and the camera casing and enclosures were designed and printed by the author (Appendix~\ref{app:mechanical}). Inherited work is identified where it is used: the vehicle dashboard and its team, and a teammate's BMS parser used only as a cross-check (Section~\ref{sec:bms-interface}). The vehicle was driven on the road sessions by Yusuf, a team member, who was not measured. The use of AI tools is declared in Section~\ref{app:ai}.

\subsection*{GA9: Independent learning ability}
The work required learning FreeRTOS timing, BLE GATT, the A76XX AT command set, on-device inference on the Raspberry Pi and vibration analysis, largely from primary sources: manufacturer schematics and manuals, standards, and live tests on the hardware. Assumptions were corrected against those sources, for example GPIO assignments against the schematic and the MQTT result handling against the modem manual (Section~\ref{sec:mqtt-networking}). A literature search was kept running so that the gap claims were tested rather than asserted once, and it narrowed two of them (Chapter~\ref{chap:lit}). The first road session taught what a vehicle test must record (markers, raw samples, a speed reference), and the second was designed from that lesson.

\subsection*{GA10: Professionalism and ethics (in part)}
The study was approved by the faculty's research ethics committee (EBE/03305/2026, Appendix~\ref{app:ethics}), and its conditions were met before the vehicle was driven: the full proposal went to the supervisor and the laboratory owner signed the risk assessment (Section~\ref{sec:ethics}). Because the approval rests on no human subjects, no driver or glance measurement was made, and the glance study is left to an ethics amendment rather than run informally. The cameras see the public, so data are minimised by design: frames are never stored, each alert keeps five fields, there is no face or number-plate recognition, and warnings are advisory. When bench records were found to carry the unit's position to a public broker, the position was removed from the published record from v0.7.5, which both nodes now run (as v0.7.8). Media that show other people are withheld, and those used are blurred (Section~\ref{sec:data-availability}). The vacation-work team's broker credentials were never given to an AI tool, as their documentation requires. Sessions with the 72\,V system energised need a second person, rated stands and chocks, and the installation is removable without drilling (R13). Results are reported honestly: what was not measured is stated as such, and errors found in the record, such as the duplicated road records, were corrected and logged.

\section{Project repository}

Firmware, camera software, analysis, planning documents and supporting tools are at \url{https://github.com/0geder/Tukzie-HPEC-HUD}; the data behind the report are indexed in \texttt{DATA\_INDEX.md} there (Section~\ref{sec:data-availability}).

\section{Declaration of Artificial Intelligence Use}
\label{app:ai}

An AI coding agent, Claude Code (Anthropic), was used throughout the project, not only for writing. It was used to draft and debug firmware, Raspberry Pi and dashboard code and their host tests; to run commands on the bench computers during tests at the author's direction; to run the automated literature search (OpenAlex, arXiv and Crossref, with every DOI re-resolved); to write the analysis package; to review the report against the supervisor's guidance; and to draft, restructure and edit report text. Decisions were made by the author, and the project log records who decided what. The author is responsible for the technical content, the analysis, the interpretations and every claim in this report.

The prompts were written by the author, and the longer ones were kept in the repository under \texttt{Planning/*-prompt.md}. Four short examples show their form, the last one of report writing:
\begin{itemize}
\item \texttt{Planning/\allowbreak{}simon-guidance-review-prompt.md} (21 October): ``Run read-only: the reviewer changes no file. Findings come back as a report; edits are decided afterwards.'' and ``Every guidance item must quote the source verbatim, with file and page''. The output was the list of additions in \texttt{Planning/\allowbreak{}winberg-guidance-compliance.md}, each of which the author accepted or rejected.
\item \texttt{Planning/\allowbreak{}run2-readiness-prompt.md} (22 October), quoting the author's brief: ``Anything still outstanding is checked and analysed; if there is a better approach, propose it and we brainstorm; if not, it is tested or placed as future work.'' with the rule ``Nothing invented. Every number comes from a file in the repository or a measurement on the day.'' The outputs were the outstanding register, the run-2 test plan, the data index and the analysis package.
\item \texttt{Planning/\allowbreak{}roadtest-photo-appendix-prompt.md} (21 October): ``Crop away or blur every face of anyone else and the vehicle's number plate.'' and ``Describe only what is visible.'' The output was Appendix~\ref{app:road1}, whose images the author checked one by one.
\item Report text (23 October). The author submitted a screenshot of the IMU pod figure with ``dont I think we revised this''. The agent confirmed that the figure and text still described revision 2, rendered the printed revision 3 from its STL files, and returned a rewritten Section~\ref{sec:mech-imu} with revision 2 kept as one sentence of history and the print table corrected. The author reviewed the section before it was kept. A standing instruction of the author applied to all drafted text: no dashes as punctuation and no unnecessary bold.
\end{itemize}

Outputs were checked rather than trusted, in five ways. Code was accepted only after it built and its host tests passed (for example 2,326 raw-stream checks for v0.7.5 and 42 tests for the analysis package) and, for firmware, after a bench test on the hardware with the flash hash verified. Numbers in the report were recomputed from the logs; when the analysis package was run on the first road session it reproduced the reported figures and found that the driving 95th percentile was 3.174, which had been written as 3.1, and the report was corrected to 3.2. A review of the report on 21 October found that the recorder had stored each road record several times (10,287 replies holding 3,464 distinct records), which had inflated the driving time from about 16 to 19.5 minutes; every affected figure was recomputed and corrected. The author corrected the agent where its reading of the hardware was wrong, for example establishing that the missing GNSS fix of 19 October came from no antenna having been fitted, not a receiver fault, and that IMU2 had not failed but had been moved to the front node (\texttt{Planning/lab-notebook.md}, 21 and 22 October). And literature claims were checked against the sources, with the gap ledger recording two claims the search narrowed.
